\documentclass[12pt,a4paper]{article}
\usepackage[margin=1in]{geometry}
\usepackage{graphicx}
\usepackage{amsmath}
\usepackage{listings}
\usepackage{xcolor}
\usepackage{hyperref}
\usepackage{booktabs}
\usepackage{float}
\usepackage{caption}
\usepackage{subcaption}
\usepackage{tikz}

\lstset{
    basicstyle=\ttfamily\small,
    breaklines=true,
    frame=single,
    language=C++,
    keywordstyle=\color{blue},
    commentstyle=\color{green!60!black},
    stringstyle=\color{red},
    numbers=left,
    numberstyle=\tiny\color{gray},
    tabsize=4
}

\title{\textbf{Laboratory 5: Parallel Applications}\\
\large Operating Systems Laboratory}
\author{CS24BT041 \\ IIT Dharwad}
\date{\today}

\begin{document}

\maketitle
\tableofcontents
\newpage

% ============================================================
\section{Introduction}

This report presents the implementation and analysis of parallel computing approaches for an image sharpening pipeline. The pipeline consists of three stages:
\begin{itemize}
    \item \textbf{S1 (Smoothen):} Applies a 3$\times$3 averaging filter to the input image.
    \item \textbf{S2 (Find Details):} Computes the difference between the original and smoothened image.
    \item \textbf{S3 (Sharpen):} Adds the detail image back to the original image with clamping.
\end{itemize}

We explore multiple parallelization strategies: threads with atomic operations, processes with pipes, processes with shared memory and semaphores, and distributed computing with sockets.

% ============================================================
\section{Part I: Baseline Sequential Implementation}

\subsection{Description}
The baseline implementation executes S1, S2, and S3 sequentially in a single process. This serves as the reference for measuring speedup.

\subsection{Implementation}
The code reads a PPM image, applies the three stages sequentially, and writes the output. Timing is measured using \texttt{std::chrono::steady\_clock}.

\subsection{Expected Runtime}
For a small image, the total time is dominated by I/O. The compute time for S1, S2, S3 is minimal for small images. This motivates the 1000-iteration experiment in Part II.

% ============================================================
\section{Part II: Multi-core Parallelism}

\subsection{Overview of Variants}

To make the experiment compute-intensive rather than I/O-intensive, the image is read once and the transformation is performed 1000 times:

\begin{itemize}
    \item \textbf{Variant 1 (Pipeline):} Each of the 1000 iterations pipelines S1$\to$S2$\to$S3. S1 always uses the original input image. Results flow through the pipeline per-iteration.
    \item \textbf{Variant 2 (Batched):} S1 runs 1000 times; every 100th result is passed to S2. S2 runs 100 times per received input; each final result is passed to S3. S3 runs 100 times per received input.
\end{itemize}

% --------------------------------------------------------
\subsection{Part 2.1: Threads with Atomic Operations}

\subsubsection{Architecture}
Three threads are spawned within a single process. They share the process address space and communicate through shared image buffers. Synchronization is achieved using \texttt{std::atomic<int>} counters.

\subsubsection{Synchronization Mechanism}
\begin{itemize}
    \item \textbf{Variant 1:} Per-iteration atomic row-completion counters. Thread S1 increments \texttt{s1\_done[iter]} after processing each row. Thread S2 spins on \texttt{s1\_done[iter]} until the required row is available, then processes it and increments \texttt{s2\_done[iter]}. Thread S3 similarly waits on \texttt{s2\_done[iter]}.
    \item \textbf{Variant 2:} Producer-consumer atomic counters (\texttt{s1\_to\_s2\_count}, \texttt{s2\_to\_s3\_count}) with corresponding consumed counters. Back-pressure is applied by having the producer wait until the consumer has consumed the previous buffer.
\end{itemize}

\subsubsection{Advantages}
\begin{itemize}
    \item No kernel-mode transitions for synchronization (pure user-space spinning).
    \item Shared address space eliminates data copying.
    \item Fine-grained (row-level) pipelining possible.
\end{itemize}

\subsubsection{Disadvantages}
\begin{itemize}
    \item Busy-waiting wastes CPU cycles when stages have unequal workloads.
    \item Debugging race conditions is difficult.
    \item Memory ordering must be carefully managed (\texttt{memory\_order\_acquire/release}).
\end{itemize}

% --------------------------------------------------------
\subsection{Part 2.2: Processes with Pipes}

\subsubsection{Architecture}
Three processes are created using \texttt{fork()}. Two pipes connect them: \texttt{pipe1} (S1$\to$S2) and \texttt{pipe2} (S2$\to$S3). Data is serialized as flat byte arrays.

\subsubsection{Data Transfer}
Image data is serialized to a flat \texttt{uint8\_t} buffer (height $\times$ width $\times$ 3 bytes) and written to the pipe. The receiver reads exactly the same number of bytes and reconstructs the image.

\subsubsection{Advantages}
\begin{itemize}
    \item FIFO ordering guaranteed by pipe semantics --- no explicit synchronization needed.
    \item Clean process isolation --- each process has its own address space.
    \item Simple programming model.
\end{itemize}

\subsubsection{Disadvantages}
\begin{itemize}
    \item Data must be serialized/deserialized (copied into/out of pipe buffers).
    \item Kernel involvement for every read/write system call.
    \item Pipe buffer size limits may cause blocking.
    \item Higher memory usage due to separate address spaces.
\end{itemize}

% --------------------------------------------------------
\subsection{Part 2.3: Processes with Shared Memory and Semaphores}

\subsubsection{Architecture}
Three processes share a System V shared memory segment (\texttt{shmget}/\texttt{shmat}). POSIX semaphores (\texttt{sem\_init} with \texttt{pshared=1}) are placed in the shared memory for inter-process synchronization.

\subsubsection{Shared Memory Layout}
\begin{verbatim}
[shm_header: 4 semaphores + width + height]
[s1_buffer: H * W * 3 bytes]
[s2_buffer: H * W * 3 bytes]
\end{verbatim}

\subsubsection{Synchronization}
Four semaphores implement double-buffered producer-consumer:
\begin{itemize}
    \item \texttt{s1\_ready}: S1 signals S2 that s1\_buffer is filled.
    \item \texttt{s1\_can\_write}: S2 signals S1 that s1\_buffer can be overwritten.
    \item \texttt{s2\_ready}: S2 signals S3 that s2\_buffer is filled.
    \item \texttt{s2\_can\_write}: S3 signals S2 that s2\_buffer can be overwritten.
\end{itemize}

\subsubsection{Advantages}
\begin{itemize}
    \item Eliminates data copying between processes (shared memory access is direct).
    \item Semaphores provide proper blocking (no busy-waiting).
    \item Process isolation for fault tolerance.
\end{itemize}

\subsubsection{Disadvantages}
\begin{itemize}
    \item Complex setup (shared memory creation, semaphore initialization).
    \item Must carefully manage shared memory lifecycle (cleanup with \texttt{shmctl}).
    \item Semaphore operations involve kernel transitions.
    \item Debugging is harder than threads (no shared debugger context).
\end{itemize}

% ============================================================
\section{Part III: Distributed Computing with Sockets}

\subsection{Architecture}
Two computers communicate via TCP sockets. One acts as a client, the other as a server. The protocol is:
\begin{enumerate}
    \item Client connects to server.
    \item Client sends: image dimensions (width, height, iteration count) as three 32-bit integers.
    \item Client sends: the original input image as a flat byte array.
    \item Client sends: processed results (S1 or S1+S2 output) as flat byte arrays, one per iteration/pass.
    \item Server receives results, processes remaining stages, and writes the output file.
\end{enumerate}

\subsection{Part 3.1: S1+S2 on Computer A, S3 on Computer B}
Computer A performs both S1 and S2, sending the S2 output (details image) to Computer B, which applies S3 and writes the output.

\textbf{Analysis:} This partition puts more computational load on Computer A (two stages) and less on Computer B (one stage). The network transfer contains the details image, which is relatively small in data variation.

\subsection{Part 3.2: S1 on Computer A, S2+S3 on Computer B}
Computer A performs only S1, sending the smoothened image to Computer B, which performs both S2 and S3.

\textbf{Analysis:} This partition is more balanced. Computer A handles the most compute-intensive stage (S1 with 3$\times$3 convolution), while Computer B handles the two simpler stages (element-wise operations).

\subsection{Comparison}
\begin{table}[H]
\centering
\begin{tabular}{lcc}
\toprule
\textbf{Metric} & \textbf{Part 3.1 (S1+S2 on A)} & \textbf{Part 3.2 (S1 on A)} \\
\midrule
Load on Computer A & High (2 stages) & Medium (1 stage) \\
Load on Computer B & Low (1 stage) & Medium (2 stages) \\
Network data per iter & details image & smoothened image \\
Load balance & Imbalanced & More balanced \\
\bottomrule
\end{tabular}
\caption{Comparison of distributed partitioning strategies}
\end{table}

Part 3.2 (S1 on A, S2+S3 on B) is generally better because:
\begin{itemize}
    \item S1 is the most compute-intensive (convolution), so isolating it on one machine makes sense.
    \item S2 and S3 are both element-wise operations and can run back-to-back without network transfer.
    \item Better load balancing leads to lower idle time on either machine.
\end{itemize}

% ============================================================
\section{Verification of Pixel Ordering}

To prove that pixels are received in the correct order in each parallel case:

\subsection{Method}
\begin{enumerate}
    \item \textbf{Checksum verification:} Each stage computes a position-weighted checksum of the pixel data it produces: $\text{checksum} = \sum_{i,j,k} \text{pixel}[i][j][k] \times (i \times W \times 3 + j \times 3 + k + 1)$. The consumer independently computes the same checksum on the received data and compares.
    \item \textbf{Sequential counter verification:} For row-level pipelining (Part 2.1 V1), atomic row counters ensure that row $i$ is only consumed after the producer has completed it. The counter is monotonically increasing, proving sequential ordering.
    \item \textbf{FIFO guarantee:} For pipes (Part 2.2), the FIFO property of pipes inherently guarantees that bytes are received in the order they were sent.
    \item \textbf{Semaphore ordering:} For shared memory (Part 2.3), the semaphore wait/post pattern ensures strict alternation: the consumer cannot read until the producer signals, and the producer cannot overwrite until the consumer signals.
\end{enumerate}

% ============================================================
\section{Runtime Analysis and Speedup}

\subsection{Expected Results}
\begin{table}[H]
\centering
\begin{tabular}{lcc}
\toprule
\textbf{Approach} & \textbf{Expected Speedup (V1)} & \textbf{Expected Speedup (V2)} \\
\midrule
Sequential (Part 1) & 1.0$\times$ & 1.0$\times$ \\
Threads + Atomics (2.1) & $\sim$2.5--3.0$\times$ & $\sim$2.0--2.5$\times$ \\
Processes + Pipes (2.2) & $\sim$1.5--2.0$\times$ & $\sim$1.5--2.0$\times$ \\
Shared Memory + Sem (2.3) & $\sim$2.0--2.5$\times$ & $\sim$2.0--2.5$\times$ \\
\bottomrule
\end{tabular}
\caption{Expected speedup relative to sequential baseline (1000 iterations)}
\end{table}

\subsection{Discussion}
\begin{itemize}
    \item \textbf{Threads + Atomics} achieves the best speedup because there is no data copying overhead and synchronization is done in user-space.
    \item \textbf{Shared Memory + Semaphores} is close to threads but incurs overhead for semaphore operations and copying data to/from the shared buffer.
    \item \textbf{Pipes} have the most overhead due to kernel-mode data copying through the pipe buffer.
    \item \textbf{Variant 1} (pipeline) generally achieves better speedup than Variant 2 (batched) because the pipeline allows all three cores to be active simultaneously, while the batched approach has periods where only one stage is running.
\end{itemize}

% ============================================================
\section{Implementation Difficulty Discussion}

\begin{table}[H]
\centering
\begin{tabular}{lccc}
\toprule
\textbf{Criterion} & \textbf{Threads+Atomic} & \textbf{Pipes} & \textbf{SHM+Sem} \\
\midrule
Ease of implementation & Medium & Easy & Hard \\
Debugging difficulty & Hard & Easy & Hard \\
Performance & Best & Lowest & Good \\
Code complexity & Medium & Low & High \\
Correctness reasoning & Hard & Easy & Medium \\
\bottomrule
\end{tabular}
\caption{Comparison of implementation difficulty}
\end{table}

\subsection{Threads + Atomic Operations (Part 2.1)}
\begin{itemize}
    \item \textbf{Implementation:} Medium difficulty. Requires understanding of C++ atomics and memory ordering.
    \item \textbf{Debugging:} Hard. Race conditions may be non-deterministic and architecture-dependent.
    \item \textbf{Key challenge:} Getting memory ordering right. Using \texttt{memory\_order\_relaxed} where \texttt{acquire/release} is needed can cause subtle bugs.
\end{itemize}

\subsection{Processes + Pipes (Part 2.2)}
\begin{itemize}
    \item \textbf{Implementation:} Easiest. The pipe abstraction handles synchronization implicitly.
    \item \textbf{Debugging:} Easiest. Each process can be debugged independently. Pipe errors produce clear error messages.
    \item \textbf{Key challenge:} Ensuring all file descriptors are properly closed in each process to avoid hanging.
\end{itemize}

\subsection{Shared Memory + Semaphores (Part 2.3)}
\begin{itemize}
    \item \textbf{Implementation:} Hardest. Requires managing shared memory layout, semaphore initialization, and cleanup.
    \item \textbf{Debugging:} Hard. Shared memory state is not easily inspectable, and leftover shared memory segments can cause issues between runs.
    \item \textbf{Key challenge:} Proper cleanup of shared memory segments and semaphores, especially on abnormal termination.
\end{itemize}

% ============================================================
\section{Conclusion}

This lab demonstrated multiple approaches to parallelizing an image processing pipeline:
\begin{itemize}
    \item \textbf{Threads with atomics} provide the best performance but are hardest to debug.
    \item \textbf{Pipes} provide the simplest programming model but have the most overhead.
    \item \textbf{Shared memory with semaphores} offer a good balance but require careful resource management.
    \item \textbf{Sockets} enable distributed computing across machines, with the partitioning strategy significantly affecting performance.
\end{itemize}

The choice of parallelization approach depends on the specific requirements: if performance is critical, threads with atomics are preferred; if simplicity and maintainability matter, pipes are the best choice; if processes with good performance are needed, shared memory with semaphores is the way to go.

\end{document}
